\documentclass[10pt,a4paper]{amsart}
\usepackage[margin=19mm]{geometry}
\usepackage{amsmath,amssymb,mathtools}
\usepackage[scr=rsfs,cal=euler]{mathalfa}
\usepackage{xfrac}
\usepackage[unicode,hidelinks,bookmarks=false]{hyperref}
\newcommand{\ie}{i.e.\ }
\newcommand{\cf}{cf.\ }
\newcommand{\resp}{respectively\ }
\newcommand{\Subsection}{\subsection}
\providecommand{\nobreakdash}{\nobreak\hbox{-}}
\setlength{\emergencystretch}{2em}
\multlinegap0pt
\usepackage{float}
\usepackage{mathtools}
\usepackage{amssymb}
\usepackage{tikz}
\usepackage[shortlabels]{enumitem}
\newtheorem{proposition}{Proposition}
\newcommand{\R}{\mathbb{R}}
\DeclareMathOperator{\Ric}{Ric}
\newcommand{\Span}[1]{\operatorname{Span}\{#1\}}
\DeclareMathOperator{\ad}{ad}
\newcommand{\frg}{\mathfrak{g}}
\newcommand{\frh}{\mathfrak{h}}
\newcommand{\g}{\lie g}
\newcommand{\h}{\lie h}
\newcommand{\lie}[1]{\mathfrak{#1}}
\title{\textbf{Almost abelian pseudo-Kähler Lie algebras}}
\author{Diego Conti, Alejandro Gil-Garc\'ia}
\date{Source: arXiv:2506.22278v2 (CC BY 4.0)}
\begin{document}
\maketitle

\noindent\emph{Source and licence.} \textbf{Almost abelian pseudo-Kähler Lie algebras}. Version arXiv:2506.22278v2 (CC BY 4.0), distributed by arXiv under the Creative Commons Attribution 4.0 International licence, http://creativecommons.org/licenses/by/4.0/, as recorded in the arXiv licence metadata for this exact version and shown on the abstract page https://arxiv.org/abs/2506.22278v2. Full licence notice: \texttt{licenses/arxiv-2506.22278-CC-BY-4.0.txt}.\par\smallskip

\noindent\textit{Editorial scope.} Counted unit: the article's proposition prop:curvature\_isotropic\_h0 (source0.tex lines 522--537). The article's complete original proof of it is delivered with it (source0.tex lines 539--595, one \texttt{proof} environment of 57 lines). No mathematics is written, repaired or re-derived: the statement and the proof are the article's own lines, in the article's own notation, and no retained line is substituted or re-flowed.  Retained uncounted as the source-local context the proof needs: lines 315--317; lines 333--335; lines 339--344; lines 362--364; lines 370--377; lines 482--490; lines 493--495.  Nothing was omitted from any retained span, so the package carries no omission record.  No retained line contains a citation, so the wrapper carries no bibliography and no bibliography entry.  No retained line contains a figure environment, an image-inclusion command or drawing code, so no image is delivered and the package declares no asset dependency.  Editorial formatting only, in addition: an AMS \texttt{amsart} preamble that restates the article's own declarations and macros used by the retained lines, namely the environments \texttt{proposition} on separate counters, together with the standard packages those lines need.  The retained environments print in this document as Proposition 1 in order of appearance, whereas the article prints the same items with the numbering of its own full document.  The complete native source of the article as distributed in the arXiv e-print for this exact version is bundled unchanged as \texttt{sources/180615/source0.tex}.
Let $\g$ be an \emph{almost abelian Lie algebra}, namely admitting a codimension one abelian ideal $\frh$; suppose $\g$ has real dimension $2n$. In particular, the derived subalgebra $[\frg,\frg]$ is abelian, and $\frg$ is two-step solvable. We can choose a basis $\{e_1,\ldots,e_{2n}\}$ of $\frg$ such that \begin{equation}\label{eqn:amostabelianbasis}
    \frh=\Span{e_1,\ldots,e_{2n-1}},\quad\ad(e_{2n})\lie h\subset\lie h.
\end{equation}

Therefore, we may choose a basis $\{e_1,\dotsc, e_{2n}\}$ of $\frg$, with dual basis $\{e^1,\dotsc, e^{2n}\}$, satisfying \eqref{eqn:amostabelianbasis} and in addition \begin{equation}\label{eq:non-isotropic_pseudo-Hermitian}
    Je_{2i-1}=e_{2i}, \quad g=\sum_{i=1}^{2n}\varepsilon_i e^i\otimes e^i, \quad \lie h_1=\Span{e_1,\dotsc, e_{2n-2}}, \quad \lie h_0=\Span{e_{2n-1},e_{2n}},
\end{equation} where $\varepsilon_i:=g(e_i,e_i)\in\{-1,+1\}$. Thus, $J\frh^\perp=\Span{e_{2n-1}}$ and $\frh^\perp=\Span{e_{2n}}$. Denote $J_1:=J|_{\frh_1}$ and $g_1:=g|_{\frh_1\times\frh_1}$.

    \begin{equation}\label{eq:D_positive_g0}
        D=\begin{pmatrix}
            A&0\\
            0&a
        \end{pmatrix}
    \end{equation} with $A\in\mathfrak{u}(\frh_1,J_1,g_1)$ and $a\in\R$.

Necessarily, $e_{2n}$ is not in $\lie h$. Defining $e_2=Je_1$ and $e_{2n-1}=-Je_{2n}$, we obtain a basis $\{e_1,\dotsc, e_{2n}\}$ of $\frg$ such that \begin{equation}\label{eqn:almosthermitianstructurezero}
    Je_{2i-1}=e_{2i}, \quad g=e^1\odot e^{2n}-e^2\odot e^{2n-1}+\sum_{i=3}^{2n-2}\varepsilon_i e^i\otimes e^i, \quad \lie h=\Span{e_1,\dotsc, e_{2n-1}},
\end{equation} where $e^i\odot e^j:=e^i\otimes e^j+e^j\otimes e^i$. In this case we consider the decomposition $$\frg=\frh^\perp\oplus J\frh^\perp\oplus V\oplus\Span{e_{2n-1}}\oplus\Span{e_{2n}},$$ where $$\frh^\perp=\Span{e_1},\quad J\frh^\perp=\Span{e_2},\quad V=\Span{e_3,\ldots,e_{2n-2}}.$$

    \begin{equation}\label{eq:D_zero_g0}
        D=\begin{pmatrix}
            -a&0&-(J_Vv)^{\flat_V}&c_1\\
            0&-a&v^{\flat_V}&c_2\\
            0&0&A&v\\
            0&0&0&a
        \end{pmatrix}
    \end{equation} with $A\in\mathfrak{u}(V,J_V,g_V)$, $v\in V$ and $a,c_1,c_2\in\R$.

\begin{proposition}\label{prop:curvature_definite_g0}
    Let $(\frg,J,g)$ be an almost abelian pseudo-Kähler Lie algebra with $\frh_0$ non-isotropic, so that $\g=\h\rtimes_D\R$ with $D$ as in \eqref{eq:D_positive_g0} and $(J,g)$ takes the form \eqref{eq:non-isotropic_pseudo-Hermitian}. Then the Levi-Civita connection is given by \begin{enumerate}
        \item $\nabla_X=0$ for $X\in\frh_1$.
        \item $\nabla_{e_{2n-1}}=a(e^{2n-1}\otimes e_{2n}-e^{2n}\otimes e_{2n-1})$.
        \item $\nabla_{e_{2n}}X=AX$ for $X\in\frh_1$ and $\nabla_{e_{2n}}e_{2n-1}=0=\nabla_{e_{2n}}e_{2n}$.
    \end{enumerate}

    Moreover, the only nonzero curvature term is given by $$R(e_{2n-1},e_{2n})=a^2(e^{2n-1}\otimes e_{2n}-e^{2n}\otimes e_{2n-1})=a\nabla_{e_{2n-1}}$$ and the Ricci curvature is given by $$\Ric=-a^2(e^{2n-1}\otimes e^{2n-1}+e^{2n}\otimes e^{2n}).$$
\end{proposition}

    The Levi-Civita connection $\nabla$ of the metric $g$ on $\frg$ is given by the Koszul formula: \begin{equation}\label{eq:Koszul}
        2g(\nabla_XY,Z)=g([X,Y],Z)-g([Y,Z],X)+g([Z,X],Y)
    \end{equation} for $X,Y,Z\in\frg$. If $X,Y,Z\in\frh$ then $g(\nabla_XY,Z)=0$. Now we compute $$2g(\nabla_XY,e_{2n})=-g([Y,e_{2n}],X)+g([e_{2n},X],Y)=g(DY,X)+g(DX,Y).$$

\begin{proposition}\label{prop:curvature_isotropic_h0}
    Let $(\frg,J,g)$ be an almost abelian pseudo-Kähler Lie algebra with $\frh_0$ isotropic, so that $\g=\h\rtimes_D\R$ with $D$ as in \eqref{eq:D_zero_g0} and $(J,g)$ takes the form \eqref{eqn:almosthermitianstructurezero}. Then the Levi-Civita connection is given by \begin{enumerate}
        \item $\nabla_{e_1}=\nabla_{e_2}=0$.
        \item $\nabla_X=0$ for $X\in V$.
        \item $\nabla_{e_{2n-1}}=-c_2(e^{2n-1}\otimes e_1+e^{2n}\otimes e_2)$.
        \item $\nabla_{e_{2n}}X=g(v,JX)e_1+g(v,X)e_2+AX$ for $X\in V$ and $$\begin{aligned}
            \nabla_{e_{2n}}e_1&=-ae_1,\\
            \nabla_{e_{2n}}e_2&=-ae_2,
        \end{aligned}\qquad\begin{aligned}
            \nabla_{e_{2n}}e_{2n-1}&=c_1e_1+v+ae_{2n-1},\\
            \nabla_{e_{2n}}e_{2n}&=c_1e_2+Jv+ae_{2n}.
        \end{aligned}$$
    \end{enumerate}

    Moreover, the only nonzero curvature term is given by $$R(e_{2n-1},e_{2n})=-3ac_2(e^{2n-1}\otimes e_1+e^{2n}\otimes e_2)=3a\nabla_{e_{2n-1}}$$ and the metric $g$ is Ricci-flat.
\end{proposition}

\begin{proof}
    The proof is analogous to that of Proposition \ref{prop:curvature_definite_g0}. We use the Koszul formula \eqref{eq:Koszul} to compute the covariant derivative operators in each of the four cases listed in the statement. We first note that $g(\nabla_XY,Z)=0$ for all $X,Y,Z\in\frh$.\medskip

    (1) We have $g(\nabla_{e_1}X,Y)=0$ for all $X,Y\in\frh$ and $g(\nabla_{e_1}e_{2n},X)=-g(e_{2n},\nabla_{e_1}X)$ for all $X\in\frg$. In particular $g(\nabla_{e_1}e_{2n},e_{2n})=0$. Therefore, we only need to compute the following terms: \begin{itemize}
        \item $2g(\nabla_{e_1}e_1,e_{2n})=-g([e_1,e_{2n}],e_1)+g([e_{2n},e_1],e_1)=2g(De_1,e_1)=-2ag(e_1,e_1)=0$ since $e_1$ is null.
        \item $2g(\nabla_{e_1}e_2,e_{2n})=-g([e_2,e_{2n}],e_1)+g([e_{2n},e_1],e_2)=g(De_2,e_1)+g(De_1,e_2)=-2ag(e_1,e_2)=0$ since $e_1$ and $e_2$ are orthogonal.
        \item Let $X\in V$, then $2g(\nabla_{e_1}X,e_{2n})=-g([X,e_{2n}],e_1)+g([e_{2n},e_1],X)=g(DX,e_1)+g(De_1,X)=0$ since $DX\in\frh$ and $\frh$ is orthogonal to $e_1$.
        \item $2g(\nabla_{e_1}e_{2n-1},e_{2n})=-g([e_{2n-1},e_{2n}],e_1)+g([e_{2n},e_1],e_{2n-1})=g(De_{2n-1},e_1)+g(De_1,e_{2n-1})=0$ since $De_{2n-1}\in\frh$ and $\frh$ is orthogonal to $e_1$.
    \end{itemize}

    We thus conclude that $\nabla_{e_1}=0$. To compute $\nabla_{e_2}$ we proceed in a similar way. The only difference is that $g(e_2,e_{2n-2})=-1$. Nevertheless, $$2g(\nabla_{e_2}e_{2n-1},e_{2n})=g(De_{2n-1},e_2)+g(De_2,e_{2n-1})=ag(e_{2n-1},e_2)+g(-ae_2,e_{2n-1})=0.$$

    Then $\nabla_{e_2}=0$ as well.\medskip

    (2) Let $X\in V$. Then for all $Y\in V$ we have $$2g(\nabla_XY,e_{2n})=g(DY,X)+g(DX,Y)=g(AY,X)+g(AX,Y)=g((A+A^*)X,Y)=0,$$ where we have used once again that $e_1$ and $e_2$ are orthogonal to $V$. We also have \begin{align*}
        2g(\nabla_Xe_{2n-1},e_{2n})&=g(De_{2n-1},X)+g(DX,e_{2n-1})=g(v,X)+g(v^tg_VXe_2,e_{2n-1})\\
        &=g(v,X)+g(v,X)g(e_2,e_{2n-1})=0,
    \end{align*} since $g(e_2,e_{2n-1})=-1$. By similar arguments to the previous case we conclude that $\nabla_X=0$ for $X\in V$.\medskip

    (3) First note that $\nabla_{e_{2n-1}}X=\nabla_Xe_{2n-1}=0$ for all $X\in\frh^\perp\oplus J\frh^\perp\oplus V$. Now we compute $$g(\nabla_{e_{2n-1}}e_{2n-1},e_{2n})=g(De_{2n-1},e_{2n-1})=g(c_2e_2,e_{2n-1})=-c_2$$ and $g(\nabla_{e_{2n-1}}e_{2n},e_{2n-1})=-g(e_{2n},\nabla_{e_{2n-1}}e_{2n-1})=c_2$. All the other terms are zero, thus we conclude that $\nabla_{e_{2n-1}}=-c_2e^{2n-1}\otimes e_1-c_2e^{2n}\otimes e_2$.\medskip

    (4) Here we make use of the fact that $\nabla$ is torsion-free: \begin{itemize}
        \item $\nabla_{e_{2n}}e_1=\nabla_{e_1}e_{2n}+[e_{2n},e_1]=De_1=-ae_1$.
        \item $\nabla_{e_{2n}}e_2=\nabla_{e_2}e_{2n}+[e_{2n},e_2]=De_2=-ae_2$.
        \item $\nabla_{e_{2n}}X=\nabla_Xe_{2n}+[e_{2n},X]=DX=v^tg_VJ_VXe_1+v^tg_VXe_2+AX=g(v,JX)e_1+g(v,X)e_2+AX$ for all $X\in V$.
        \item $\nabla_{e_{2n}}e_{2n-1}=\nabla_{e_{2n-1}}e_{2n}+[e_{2n},e_{2n-1}]=-c_2e_2+De_{2n-1}=c_1e_1+v+ae_{2n-1}$.
    \end{itemize}

    Finally we compute $\nabla_{e_{2n}}e_{2n}$: \begin{itemize}
        \item $g(\nabla_{e_{2n}}e_{2n},e_1)=-g(e_{2n},\nabla_{e_{2n}}e_1)=-g(e_{2n},-ae_1)=a$.
        \item $g(\nabla_{e_{2n}}e_{2n},e_2)=-g(e_{2n},\nabla_{e_{2n}}e_2)=-g(e_{2n},-ae_2)=0$.
        \item $g(\nabla_{e_{2n}}e_{2n},X)=-g(e_{2n},\nabla_{e_{2n}}X)=-g(e_{2n},g(v,JX)e_1)=-g(v,JX)=g(Jv,X)$.
        \item $g(\nabla_{e_{2n}}e_{2n},e_{2n-1})=-g(e_{2n},\nabla_{e_{2n}}e_{2n-1})=-g(e_{2n},c_1e_1)=-c_1$.
    \end{itemize}

    Therefore we get $\nabla_{e_{2n}}e_{2n}=c_1e_2+Jv+ae_{2n}$.\medskip
    
    Let us compute now the Riemann curvature tensor of $g$. We immediately see that $R(X,Y)=0$ for all $X,Y\in\frh$. Now we have $$R(e_1,e_{2n})=-\nabla_{[e_1,e_{2n}]}=\nabla_{De_1}=-a\nabla_{e_1}=0,$$ and similarly $R(e_2,e_{2n})=R(X,e_{2n})=0$ for all $X\in V$. Using that $\nabla_{De_{2n-1}}=a\nabla_{e_{2n-1}}$ we get $$R(e_{2n-1},e_{2n})=\nabla_{e_{2n-1}}\nabla_{e_{2n}}-\nabla_{e_{2n}}\nabla_{e_{2n-1}}+a\nabla_{e_{2n-1}}.$$

    The above expression is not identically zero only when evaluated in $e_{2n-1}$ and $e_{2n}$. Indeed: \begin{align*}
        R(e_{2n-1},e_{2n})e_{2n-1}&=\nabla_{e_{2n-1}}\nabla_{e_{2n}}e_{2n-1}-\nabla_{e_{2n}}\nabla_{e_{2n-1}}e_{2n-1}+a\nabla_{e_{2n-1}}e_{2n-1}\\
        &=\nabla_{e_{2n-1}}(c_1e_1+v+ae_{2n-1})-\nabla_{e_{2n}}(-c_2e_1)+a(-c_2e_1)\\
        &=a\nabla_{e_{2n-1}}e_{2n-1}+c_2\nabla_{e_{2n}}e_1-ac_2e_1\\
        &=-3ac_2e_1,\\
        R(e_{2n-1},e_{2n})e_{2n}&=R(e_{2n-1},e_{2n})Je_{2n-1}=JR(e_{2n-1},e_{2n})e_{2n-1}\\
        &=-3ac_2e_2.
    \end{align*}
    Therefore we conclude that $R(e_{2n-1},e_{2n})=-3ac_2(e^{2n-1}\otimes e_1+e^{2n}\otimes e_2)=3a\nabla_{e_{2n-1}}$.\medskip

    Let $E_1,\ldots,E_{2n}$ be an orthonormal basis of $\frg$, where $E_i=e_i$ for $i=3,\ldots,2n-2$ and $$E_1=\tfrac{1}{\sqrt{2}}(e_1+e_{2n}),\quad E_2=\tfrac{1}{\sqrt{2}}(e_2+e_{2n-1}),\quad E_{2n-1}=\tfrac{1}{\sqrt{2}}(e_2-e_{2n-1}),\quad E_{2n}=\tfrac{1}{\sqrt{2}}(e_1-e_{2n}).$$

    Since $R(X,w)=0$ for all $X\in\frh^\perp\oplus J\frh^\perp\oplus V$, $\Ric(X,w)=\sum_{i=1}^{2n}\varepsilon_ig(R(E_i,X)w,E_i)=0$. Moreover, since pseudo-K\"ahler metrics satisfy $\Ric(X,JX)=0$, we also have that $\Ric(e_{2n-1},e_{2n})=\Ric(e_{2n-1},Je_{2n-1})=0$. Now we compute \begin{align*}
        \Ric(e_{2n-1},e_{2n-1})&=\sum_{i=1}^{2n}\varepsilon_ig(R(E_i,e_{2n-1})e_{2n-1},E_i)\\
        &=\varepsilon_1g(R(E_1,e_{2n-1})e_{2n-1},E_1)+\varepsilon_{2n}g(R(E_{2n},e_{2n-1})e_{2n-1},E_{2n}),
    \end{align*} where $\varepsilon_1=g(E_1,E_1)=1$ and $\varepsilon_{2n}=g(E_{2n},E_{2n})=-1$. The first term gives $\frac{3}{2}ac_2$ and the second gives $-\frac{3}{2}ac_2$. Therefore $\Ric(e_{2n-1},e_{2n-1})=0$. Finally,
    \[\Ric(e_{2n},e_{2n})=\Ric(Je_{2n-1},Je_{2n-1})=\Ric(e_{2n-1},e_{2n-1})=0.\] This concludes the proof that the metric $g$ is Ricci-flat.
\end{proof}

\end{document}
